% !TEX program = lualatex
\documentclass[11pt,a4paper]{article}

\usepackage[english]{babel}
\usepackage{tensygrid-hackathon}
\usepackage{booktabs,tabularx,enumitem}

\newcommand{\eventdate}{26 November 2026}
\newcommand{\eventtime}{09:00--17:00}
\newcommand{\eventvenue}{Sala InnoEnergy, ETSEIB, Barcelona}
\newcommand{\eventlanguage}{English}
\newcommand{\eventversion}{v1.2, 26 November 2026}

\begin{document}

\pagenumbering{gobble}

\HackathonHeader{Keep the grid stable}
{Dynamic topology optimization \& converter tuning}
{\textbf{TenSyGrid 1-day hackathon challenge}\\[1pt]
{\footnotesize \color{tggrey}Universitat Polit\`ecnica de Catalunya \textperiodcentered\ CETPartnership Horizon Europe \mbox{(GA 101069750)}}}

% --- INTRO TEXT ---
\small
Power systems are rapidly transitioning towards inverter-dominated operation. This 1-day hackathon challenges teams
to first \textbf{reconstruct a function from structural and coefficient matrices} in an ungraded one-hour warm-up, then use the VeraGrid Python API to \textbf{minimize transmission losses} via topology switching while \textbf{preserving (and improving)
  small-signal stability} through converter control tuning. Submissions are scored automatically across multiple operating scenarios.

\vspace{6pt}

% --- EVENT INFO ---
\normalsize
\begin{minipage}[t]{0.48\textwidth}
  \renewcommand{\arraystretch}{1.25}
  \begin{tabularx}{\textwidth}{@{}lX@{}}
    \textbf{\color{tgblue}Date} & \eventdate \\
    \textbf{\color{tgblue}Time} & \eventtime \\
  \end{tabularx}
\end{minipage}
\hfill
\begin{minipage}[t]{0.48\textwidth}
  \renewcommand{\arraystretch}{1.25}
  \begin{tabularx}{\textwidth}{@{}lX@{}}
    \textbf{\color{tgblue}Venue}    & \eventvenue    \\
    \textbf{\color{tgblue}Language} & \eventlanguage \\
  \end{tabularx}
\end{minipage}

\vspace{10pt}
{\large\bfseries\color{tgblue}Programme}\par
\vspace{4pt}

\small
\renewcommand{\arraystretch}{1.10}
\begin{tabularx}{\textwidth}{@{}p{26mm}X@{}}
  \toprule
  \textbf{\color{tgblue}Time} & \textbf{\color{tgblue}Session and milestone}                                                                                                              \\
  \midrule
  \textbf{09:00--09:15}       & \textbf{Hackathon kick-off:} Challenge presentation, stability criteria, and VeraGrid platform introduction by Carlos Herrera Vazquez.                    \\
  \textbf{09:15--10:15}       & \textbf{Ungraded warm-up: Reconstruct the function.} From matrices $S$ and $\Phi$, build expressions and a Python evaluator; validate on reference cases. \\
  \textbf{10:15--11:00}       & \textbf{Part 2: Static grid and topology switching.} Baseline evaluation, topology and dispatch optimisation with VeraGrid.                               \\
  \textbf{11:00--11:30}       & \textit{\color{tgpetrol}Coffee break.}                                                                                                                    \\
  \textbf{11:30--12:00}       & \textbf{Hackathon support:} Dedicated mentoring and technical Q\&A with TenSyGrid consortium researchers.                                                 \\
  \textbf{12:00--13:00}       & \textbf{Part 3: Dynamic layer and converter control tuning.} Computing eigenvalues, modal analysis, and stabilizing critical oscillation modes.           \\
  \textbf{13:00--14:00}       & \textit{\color{tgpetrol}Lunch break.}                                                                                                                     \\
  \textbf{14:00--15:30}       & \textbf{Part 4: Automation and batch resilience.} Multi-scenario batch sweep across 300 test operating points and parameter optimization.                 \\
  \textbf{15:30--16:00}       & \textit{\color{tgpetrol}Coffee break and code freeze.} Submission of final scripts to the automated scoring harness.                                      \\
  \textbf{16:00--17:00}       & \textbf{Hackathon wrap-up and winners announcement:} Team presentations, live leaderboard reveal, and awards ceremony.                                    \\
  \bottomrule
\end{tabularx}

\vfill
\hfill{\footnotesize\color{tggrey}\textbf{\eventversion}}

\end{document}

